\SourcePage{569}{572}
\section{Isotopic invariance}

At present there is still no complete theory of the so-called \emph{nuclear
forces}: the forces acting between nuclear particles (\emph{nucleons}) and
binding them together within an atomic nucleus. Consequently, for now the
description of nuclear forces must rely on experiment to a much greater extent
than would be necessary if a consistent theory were available.

The two kinds of particle belonging to the nucleon family differ first of all
in their electrical properties: protons ($p$) carry positive charge, whereas
neutrons ($n$) are electrically neutral. Both kinds nevertheless have the same
spin $1/2$, and their masses are nearly equal (the proton and neutron masses
are respectively $1836{,}1$ and $1838{,}6$ electron masses). This likeness is
not accidental. In spite of their different electrical properties, the proton
and neutron are particles very much alike, and this likeness is of fundamental
importance.

It is found that, when the relatively weak electrical forces are set aside,
the interaction between two protons closely resembles that between two
neutrons. This property of nuclear forces is called \emph{charge
symmetry}\footnote{One of its manifestations is the similarity in the
properties (binding energy, energy spectrum, and so forth) of the so-called
mirror nuclei, namely nuclei obtained from one another by replacing every
proton by a neutron and conversely.}.

To the accuracy with which this symmetry holds, one may say in particular that
the two-proton ($pp$) and two-neutron ($nn$) systems possess states with the
same properties. It is important here, of course, that protons and neutrons
obey the same statistics (Fermi statistics). Hence only states having the same
wave-function symmetry are possible for the $pp$ and $nn$ systems: the
functions $\psi(\mathbf r_1,\sigma_1;\mathbf r_2,\sigma_2)$ must be
antisymmetric under simultaneous interchange of the particles' coordinates
and spins.

\SourcePage{570}{573}
Charge symmetry, however, is only one manifestation of a still deeper physical
likeness between proton and neutron, known as \emph{isotopic
invariance}\footnote{This invariance is also called \emph{isobaric} in the
literature.}. This deeper regularity establishes an analogy not only between
the $pp$ and $nn$ systems (which are transformed into one another by exchanging
all protons and neutrons), but also with the $pn$ system, whose particles are
different. A complete analogy is plainly impossible, since the allowed states
of a $pn$ system, being states of nonidentical particles, need not in any event
be restricted to antisymmetric wave functions. Among the possible states of
the $pn$ system, however, there are states whose properties agree to high
accuracy with those of a pair of identical nucleons\footnote{This was
established from an analysis of experimental data on the scattering of
neutrons and protons by protons by \emph{Breit, Condon, and Present}
(G.~Breit, E.~U.~Condon, R.~D.~Present, 1936).}; naturally, these states are
described by antisymmetric wave functions. The remaining $pn$ states have
symmetric wave functions and have no counterparts in the $pp$ and $nn$
systems.

Like charge symmetry, isotopic invariance applies only when electromagnetic
interaction is neglected. A second reason for its approximate character is
the small difference between neutron and proton masses: exact neutron--proton
symmetry would plainly require their masses to coincide exactly\footnote{It is
reasonable to suppose that the neutron--proton mass difference is itself
electromagnetic in origin.}.

A convenient formal apparatus can be introduced to describe isotopic
invariance. Its natural route follows from observing that isotopic invariance
amounts to the possibility of classifying the states of a nucleon system by
the symmetry of its coordinate--spin wave functions $\psi$, without regard to
which of the two kinds its nucleons belong. The required apparatus must
therefore provide a new quantum number for characterizing the system's states,
such that specifying this number uniquely fixes the symmetry of $\psi$. We
have already encountered an analogous situation in studying a system of
spin-$1/2$ particles. In fact, we saw (see \S~63) that specifying the total spin
$S$ of such a system uniquely determines the symmetry of

\SourcePage{571}{574}
its coordinate wave function $\varphi$, independently of which of the two
possible values ($\pm1/2$) are taken by the spin projections $\sigma$ of the
individual particles.

It is therefore natural, in a formal description of isotopic invariance, to
regard the neutron and proton as two different \emph{charge states} of one and
the same particle (the nucleon). They are distinguished by the projection of a
new vector $\boldsymbol\tau$, formally analogous to a spin-$1/2$ vector. This
new quantity, customarily called \emph{isotopic spin} (or simply
\emph{isospin})\footnote{It was introduced by \emph{Heisenberg} (1932) and
applied to the description of isotopic invariance by \emph{Cassen and Condon}
(B.~Cassen, E.~U.~Condon, 1936).}, is a vector in an auxiliary ``isotopic
space'' $\xi\eta\zeta$, which of course has nothing to do with physical space.

The projection of a nucleon's isotopic spin on the $\zeta$ axis can take only
the two values $\tau_\zeta=\pm1/2$. By convention, $+1/2$ is assigned to the
proton and $-1/2$ to the neutron\footnote{The opposite convention is also used
in the literature.}. The isospins of several nucleons combine into the total
isospin of the system according to the ordinary rules for adding spins. The
$\zeta$ component of the system's total isospin is then the sum of the
$\tau_\zeta$ values of all its constituent particles. For a nucleus with
proton number (that is, atomic number) $Z$, neutron number $N$, and mass number
$A=Z+N$, we have
\begin{equation}
 T_\zeta=\sum\tau_\zeta=\frac{Z-N}{2}=Z-\frac A2,
 \tag{116.1}
\end{equation}
so, for a specified number of nucleons, $T_\zeta$ determines the total charge
of the system. Thus $T_\zeta$ is strictly conserved, this being simply an
expression of charge conservation.

The magnitude $T$ of the system's total isotopic spin determines the symmetry
of the ``charge part'' $\omega$ of its wave function, just as the total spin
$S$ determines the symmetry of the spin wave function. It thereby also fixes
the symmetry of the coordinate--spin (that is, ordinary) wave function $\psi$,
since the complete wave function of the nucleon system (the product
$\psi\omega$) must have a definite symmetry. As for any fermions, it must be
antisymmetric under simultaneous interchange of the coordinates, spins, and
``charge variables'' $\tau_\zeta$ of the particles. Thus the definite symmetry
of the wave functions $\psi$

\SourcePage{572}{575}
for any nucleon system is represented in this scheme precisely by conservation
of $T$.

Equivalently, isotopic invariance means that the system's properties are
invariant under arbitrary rotations in isotopic space. States differing only
in the value of $T_\zeta$, with $T$ and all other quantum numbers fixed, have
the same properties. Charge symmetry in particular---invariance under
interchanging every neutron with a proton and conversely---is a special case
of isotopic invariance. In this description it appears as invariance under
simultaneous reversal of all $\tau_\zeta$, that is, under a rotation through
$180^\circ$ in isospace about an axis lying in the $\xi\eta$ plane.

The evident violation of isotopic invariance by the Coulomb interaction is
also formally apparent in this scheme: the Coulomb interaction depends on
charge, hence on the $\zeta$ components of isospin, which are not invariant
under rotations in $\xi\eta\zeta$ space.

Consider, for example, a system of two nucleons. Its total isotopic spin can
have the values $T=1$ and $T=0$. For $T=1$, the possible projection values are
$T_\zeta=1,0,-1$. According to (116.1), these correspond to charges $2$, $1$,
and $0$; thus a system with $T=1$ can be realized as $pp$, $pn$, or $nn$. The
charge part $\omega$ of the wave function is symmetric for $T=1$ (just as a
spin value $S=1$ corresponds to a symmetric spin function; compare \S~62).
Consequently, $T=1$ corresponds to states with antisymmetric ordinary wave
functions $\psi$. For $T=0$, only $T_\zeta=0$ is possible and the corresponding
$\omega$ is antisymmetric; this case therefore comprises the $pn$ states with
symmetric wave functions $\psi$.

Isotopic spin has an associated operator $\hat{\boldsymbol\tau}$ acting on the
charge variable $\tau_\zeta$ in the wave function in the same manner as the
spin operator $\hat{\mathbf s}$ acts on the spin variable $\sigma$. Because the
formal analogy is complete, the operators $\hat\tau_\xi$, $\hat\tau_\eta$,
$\hat\tau_\zeta$ are represented by the same Pauli matrices (55.7) as
$\hat s_x$, $\hat s_y$, $\hat s_z$.

We note several combinations of these operators that have a simple, direct
meaning. The sum
\[
 \hat\tau_+=\hat\tau_\xi+i\hat\tau_\eta=
 \begin{pmatrix}0&1\\0&0\end{pmatrix}
\]
is an operator which turns a neutron wave function into a proton wave function

\SourcePage{573}{576}
when applied to it, while giving zero on a proton function. Similarly, the
operator
\[
 \hat\tau_-=\hat\tau_\xi-i\hat\tau_\eta=
 \begin{pmatrix}0&0\\1&0\end{pmatrix}
\]
turns a proton into a neutron and annihilates a neutron. Finally, the operator
\[
 \frac12+\hat\tau_\zeta=\begin{pmatrix}1&0\\0&0\end{pmatrix}
\]
leaves a proton function unchanged and annihilates a neutron function; it may
be called the nucleon charge operator (in units of $e$).

We shall also show how the operator $\hat P$ that interchanges two particles
can be expressed through their isospin operators
$\hat{\boldsymbol\tau}_1$, $\hat{\boldsymbol\tau}_2$. By definition, its action
on the wave function $\psi(\mathbf r_1,\sigma_1;\mathbf r_2,\sigma_2)$ of a
two-particle system interchanges the coordinates and spins of the particles,
that is, the variables $\mathbf r_1,\sigma_1$ and
$\mathbf r_2,\sigma_2$. Its eigenvalues are $\pm1$, obtained when it acts on a
symmetric or an antisymmetric function $\psi$:
\begin{equation}
 \hat P\psi_{\text{sym}}=\psi_{\text{sym}},\qquad
 \hat P\psi_{\text{antisym}}=-\psi_{\text{antisym}}.
 \tag{116.2}
\end{equation}

We saw above that the functions $\psi_{\text{sym}}$ and
$\psi_{\text{antisym}}$ correspond respectively to charge functions
$\omega_T$ with total isospin $T=0$ and $T=1$. Hence, if $\hat P$ is to be
represented in a form acting on charge variables, it must have the properties
\begin{equation}
 \hat P\omega_0=\omega_0,\qquad \hat P\omega_1=-\omega_1.
 \tag{116.3}
\end{equation}
These conditions are satisfied by the operator $1-\hat{\mathbf T}^2$. This is
seen immediately by noting that $\omega_T$ is an eigenfunction of
$\hat{\mathbf T}^2$ with eigenvalue $T(T+1)$. Finally, putting
$\hat{\mathbf T}=\hat{\boldsymbol\tau}_1+\hat{\boldsymbol\tau}_2$ and using the
fact that $\hat{\boldsymbol\tau}_1^2$ and $\hat{\boldsymbol\tau}_2^2$ have the
same definite value $\tau(\tau+1)=3/4$, we obtain the required final
expression\footnote{An operator of this form, constructed from the ordinary
spins of the particles, was encountered earlier in the problems to \S~62.}
\begin{equation}
 \hat P=1-\hat{\mathbf T}^2=-\frac12
 -2\hat{\boldsymbol\tau}_1\hat{\boldsymbol\tau}_2.
 \tag{116.4}
\end{equation}

\SourcePage{574}{577}
There are definite isospin selection rules for matrix elements of various
physical quantities of a nucleon system (L.~A.~Radicati, 1952). Let $F$ be a
quantity (of arbitrary tensor character) which is additive in the sense that
its value for the system equals the sum of its values for the individual
nucleons. Write its operator as
\[
 \hat F=\sum_p\hat f_p+\sum_n\hat f_n,
\]
where the sums extend over all protons and all neutrons in the system. This
expression may identically be rewritten as
\begin{equation}
 \begin{split}
 \hat F&=\sum\left(\frac12+\hat\tau_\zeta\right)\hat f_p
 +\sum\left(\frac12-\hat\tau_\zeta\right)\hat f_n\\
 &=\frac12\sum(\hat f_p+\hat f_n)
 +\sum(\hat f_p-\hat f_n)\hat\tau_\zeta,
 \end{split}
 \tag{116.5}
\end{equation}
where every sum in the last expression runs over all nucleons, both protons
and neutrons. The first term in (116.5) is a scalar, while the second is the
$\zeta$ component of a vector in isospace. They therefore obey the same
isospin selection rules as those applying to scalars and vectors in ordinary
space with respect to orbital angular momentum (see \S~29). An isotopic scalar
allows only transitions for which $T$ does not change. The $\zeta$ component
of an isotopic vector has matrix elements only for transitions with
$\Delta T=0,\pm1$; in addition, $\Delta T=0$ transitions between states with
$T_\zeta=0$ are forbidden, that is, for systems having equal numbers of
neutrons and protons. The last rule follows because the matrix element of a
transition with $\Delta T=0$ is proportional to $T_\zeta$; see (29.7).

For the nuclear dipole moment, for example, the quantities $f_p$ are the
products $e\mathbf r$, while $f_n=0$. The first term in (116.5) is then
\[
 \frac e2\sum\mathbf r=\frac{e}{2m}\sum m\mathbf r,
\]
that is, it is proportional to the radius vector of the center of mass and can
be made zero by a suitable choice of origin. In other words, the nuclear
dipole moment reduces to the $\zeta$ component of an isotopic vector.
